The Reock score (Reock $= A / A(\text{MBC})$) compares a district's area to its
minimum bounding circle (MBC), providing a compactness measure that is substantially
less sensitive to boundary representation artifacts than the Polsby-Popper score.
We study Reock across four \textsc{bisect} structure algorithms---standard-bisect,
prime-factor, ratio-optimal, and moving-knife---for North Carolina ($k=14$),
Wisconsin ($k=8$), and Texas ($k=38$) under 2020 Census data.
We prove that Welzl's randomised algorithm computes the exact MBC in $O(n)$
expected time and is implemented in \texttt{crates/bisect-analysis/src/compactness.rs}.
Key findings: (1) Reock decreases by $\leq 0.02$ when switching from smoothed to
raw TIGER/Line boundaries on NC coastal districts, versus PP decreasing by
$0.04$--$0.08$, confirming Reock's boundary insensitivity.
(2) The moving-knife algorithm (MKA) achieves mean district Reock $\geq 0.38^\dagger$
on NC 2020 versus $0.31^\dagger$ for standard-bisect---a 23\% improvement in
the minimum Reock score.
(3) The empirical Reock--PP gap averages $+0.12$ on NC 2020 under ratio-optimal,
confirming that Reock is systematically more lenient than PP.
All results use single seed $= 42^\dagger$.
